% TeX migration of the admitted phase258 manuscript; no mathematical edits.
% Exact per-page extracted text is checked against the frozen nine-page PDF.
% Options for packages loaded elsewhere
\PassOptionsToPackage{unicode}{hyperref}
\PassOptionsToPackage{hyphens}{url}
\documentclass[
  11pt,
]{article}
\usepackage{xcolor}
\usepackage[margin=27mm]{geometry}
\usepackage{amsmath,amssymb}
\setcounter{secnumdepth}{-\maxdimen} % remove section numbering
\usepackage{iftex}
\ifPDFTeX
  \usepackage[T1]{fontenc}
  \usepackage[utf8]{inputenc}
  \usepackage{textcomp} % provide euro and other symbols
\else % if luatex or xetex
  \usepackage{unicode-math} % this also loads fontspec
  \defaultfontfeatures{Scale=MatchLowercase}
  \defaultfontfeatures[\rmfamily]{Ligatures=TeX,Scale=1}
\fi
\usepackage{lmodern}
\ifPDFTeX\else
  % xetex/luatex font selection
\fi
% Use upquote if available, for straight quotes in verbatim environments
\IfFileExists{upquote.sty}{\usepackage{upquote}}{}
\IfFileExists{microtype.sty}{% use microtype if available
  \usepackage[]{microtype}
  \UseMicrotypeSet[protrusion]{basicmath} % disable protrusion for tt fonts
}{}
\makeatletter
\@ifundefined{KOMAClassName}{% if non-KOMA class
  \IfFileExists{parskip.sty}{%
    \usepackage{parskip}
  }{% else
    \setlength{\parindent}{0pt}
    \setlength{\parskip}{6pt plus 2pt minus 1pt}}
}{% if KOMA class
  \KOMAoptions{parskip=half}}
\makeatother
\usepackage{longtable,booktabs,array}
\newcounter{none} % for unnumbered tables
\usepackage{calc} % for calculating minipage widths
% Correct order of tables after \paragraph or \subparagraph
\usepackage{etoolbox}
\makeatletter
\patchcmd\longtable{\par}{\if@noskipsec\mbox{}\fi\par}{}{}
\makeatother
% Allow footnotes in longtable head/foot
\IfFileExists{footnotehyper.sty}{\usepackage{footnotehyper}}{\usepackage{footnote}}
\makesavenoteenv{longtable}
\setlength{\emergencystretch}{3em} % prevent overfull lines
\providecommand{\tightlist}{%
  \setlength{\itemsep}{0pt}\setlength{\parskip}{0pt}}
\usepackage{bookmark}
\IfFileExists{xurl.sty}{\usepackage{xurl}}{} % add URL line breaks if available
\urlstyle{same}
% fallback for those not using the hyperref driver hyperxmp:
\makeatletter
\@ifundefined{xmpquote}{\newcommand{\xmpquote}[1]{#1}}{}
\makeatother
\hypersetup{
  pdftitle={Centralizer extensions of projective-line maniplexes},
  pdfauthor={human\_compression project --- private author draft},
  hidelinks,
  pdfcreator={LaTeX via pandoc}}

\title{Centralizer extensions of projective-line maniplexes}
\author{human\_compression project --- private author draft}
\date{10 October 2026}

\begin{document}
\maketitle

\subsection{1. Results and dependencies}\label{results-and-dependencies}

A string marking of a group \(G\) is a generating ordered tuple
\((r_0,\ldots,r_{n-1})\) of distinct nonidentity involutions such that
\(r_i r_j=r_j r_i\) whenever \(|i-j|>1\). We do not impose the
intersection property. A group with such a marking is often called an
sggi. We write \(G_I=\langle r_i:i\in I\rangle\) and
\(G_i=\langle r_j:j\ne i\rangle\). Its type is the sequence of adjacent
product orders. Put \(L(q)=\operatorname{PSL}_2(q)\) and
\(P(q)=\operatorname{PGL}_2(q)\).

This paper recovers the principal results of Leemans--Toledo {[}LT{]}
from the common projective-line theory {[}E{]}. The change of category
requires a proof: {[}E{]}'s classification assumes the intersection
property, whereas a string marking need not have it. Sections 3--4
supply this interface. We then recover the source's two infinite
constructions. A third infinite construction corrects its introduction
and concluding face-count claim. The exact edition under consideration
is the eleven-page author manuscript dated 15 May 2023, bound in the
accompanying source ledger.

\textbf{Theorem 1 (rank and local groups; Leemans--Toledo).} If \(L(q)\)
has a string marking of rank at least four, then \(q\) is odd, the rank
is four, and each end group \(G_0,G_3\) is isomorphic to \(S_4,A_5\), or
\(L(q)\). If both end groups are proper, there is exactly one marking up
to isomorphism for \(q=11\), of type \(\{3,5,3\}\), and exactly one for
\(q=19\), of type \(\{5,3,5\}\). These are the 11-cell and 57-cell; both
satisfy IP.

\textbf{Theorem 2 (extensions and the two source families).}

\begin{enumerate}
\def\labelenumi{\arabic{enumi}.}
\tightlist
\item
  Let \(q\) be odd, choose \(\alpha\in\{1,-1\}\) with
  \(q\equiv\alpha\pmod4\), and put \(N=(q-\alpha)/2\). Every rank-three
  string marking of \(L(q)\) whose first adjacent order is an odd
  divisor of \(N\) extends, with the same group, to rank four.
\item
  For every prime \(p\equiv1\pmod{12}\) there is such a marking with
  both end groups equal to \(L(p)\). Its first three generators have
  type \(\{3,p\}\).
\item
  For every prime \(p\equiv\pm1\pmod{20}\), \(p\ne19\), there is such a
  marking with one end group \(A_5\) and the other \(L(p)\).
\end{enumerate}

Both families in parts 2--3 are infinite and nonorientable. These are
{[}LT, Theorem 4.4, Lemma 4.5, Proposition 4.7 and their corollaries{]}.

\textbf{Theorem 3 (the omitted symmetric-end family).} For every prime
\(p\equiv\pm1\pmod{24}\) there is a rank-four string marking of \(L(p)\)
with facet group \(S_4\) of type \(\{3,4\}\) and vertex group \(L(p)\).
Consequently the associated regular maniplex has one vertex and
\(|L(p)|/24\) facets. There are infinitely many such examples.

Theorem 3 is a correction to the restricted face-count statement in the
source's introduction and conclusion. It is consistent with its
correctly stated Theorem 3.10 and Remark 4.1, which allow \(S_4\). No
priority claim or new open-question credit is attached to the
correction. Section 8 gives small exact witnesses for it and for two
other literal source errors. The final PGL conjecture in {[}LT, §5{]} is
treated in a separate packet; it is not silently included in the
recovery of the above known results.

We retain the full common paper {[}E{]}, including its classical Dickson
subgroup input and exact finite-amalgam program. The parts used here are
the characteristic-two involution centralizer, the centralizer descent
for a nonabelian dihedral group, the six exceptional local markings and
the ten universal presentations. We additionally retain the usual cyclic
tori and their normalizers, and the maximal \(A_5\) and \(S_4\) cases of
Dickson's theorem as printed in Giudici {[}G, Theorem 2.2(6),(8),
p.4{]}. Dirichlet's theorem supplies infinitude of primes in the stated
progressions. These outside theorems are not reproved or formalized
here.

\subsection{2. Marked groups and regular
maniplexes}\label{marked-groups-and-regular-maniplexes}

We recall the group--object correspondence so that the category change
is explicit. Given a string marking, let the flags be the elements of
\(G\) and let the \(i\)-adjacency be left multiplication by \(r_i\).
Each adjacency is fixed-point-free and involutory. The product of two
distinct adjacencies is fixed-point-free because \(r_i r_j\ne1\).
Nonconsecutive adjacencies commute, and connectedness is precisely
generation of \(G\). Thus these data satisfy the definition of a
maniplex in {[}LT, §2{]}.

A permutation of the flags preserves every adjacency exactly when it
commutes with every left multiplication by a generator, hence with the
left regular action of \(G\). Such a permutation is right multiplication
by its value at the identity: if \(f(1)=h\), then \(f(g)=gh\).
Conversely every right multiplication commutes with the left action. The
full colour-preserving automorphism group is therefore the right regular
group, isomorphic to \(G\) by sending \(h\) to multiplication on the
right by \(h^{-1}\), and the maniplex is regular. This establishes the
full automorphism group, rather than merely a flag-transitive subgroup.

Conversely, in a regular maniplex choose a base flag and the unique
automorphism taking it to each adjacent flag. Automorphisms commute with
the adjacency permutations. Moving the base flag along a coloured path
therefore expresses its terminal automorphism as the corresponding
reversed product of these distinguished automorphisms. They generate,
are distinct nonidentity involutions, and commute at nonconsecutive
positions. Changing the base flag conjugates the marking. This gives the
claimed correspondence up to marking-preserving isomorphism.

An \(i\)-face is an orbit of the subgroup on every colour except \(i\).
For our regular model its number is \([G:G_i]\). Reversing the marking
interchanges vertices and facets and gives the dual maniplex. In
particular a reversal automorphism of the marked group gives a duality.
The even-length word subgroup has index at most two. If \(G\) is a
nonabelian simple group it is all of \(G\): it is normal, and the
quotient is generated by the common image of the involutions, of order
at most two. This is exactly the source's nonorientability criterion.

\subsection{3. Centralizer bounds without the intersection
property}\label{centralizer-bounds-without-the-intersection-property}

We record precisely the two results borrowed from the umbrella. In
characteristic two the centralizer of a nonidentity involution in
\(P(q)\) is abelian {[}E, Lemma 3{]}. In odd characteristic, if \(a,b\)
are involutions and \(|ab|>2\), their common centralizer in \(P(K)\) has
at most two elements. It is trivial in the unipotent case; otherwise its
unique nonidentity element is represented by

\[ A-\tfrac12\operatorname{tr}(A)I, \tag{1} \]

where \(A\) represents \(ab\). If \(a,b\) are defined over a subfield
\(k\), this whole centralizer lies in \(P(k)\) {[}E, Theorem 5{]}.
Neither centralizer statement uses IP. Their complete proofs remain in
{[}E{]}.

\textbf{Lemma 4 (a string endpoint).} In a centre-free group with a
string marking, \(r_0\) and \(r_1\) do not commute.

\emph{Proof.} Otherwise \(r_0\) commutes with \(r_1\) and, by the string
relations, with all remaining generators. It then commutes with every
word in the generators, so is central and hence trivial, a
contradiction. \(\square\)

\textbf{Lemma 5 (rank obstruction).} Suppose a centre-free group has the
property that any two noncommuting involutions have at most one
nonidentity common centralizer element. Every string marking of that
group has rank at most four.

\emph{Proof.} At rank at least five, Lemma 4 makes \(r_0,r_1\)
noncommuting. Both \(r_3\) and \(r_4\) centralize them. The hypothesis
forces \(r_3=r_4\), contradicting distinctness of the marking.
\(\square\)

In characteristic two, a rank-four marking of \(L(q)\) is also
impossible: \(r_0,r_1\) both centralize \(r_3\), whose centralizer is
abelian. This contradicts Lemma 4. For \(q\ge4\), the group \(L(q)\) is
nonabelian simple and centre-free, so these arguments apply. For
\(q=2\), there are only three involutions in \(S_3\); for \(q=3\), the
involutions of \(A_4\) generate only its four-group. These dispose of
the small groups independently. We have proved the entire rank assertion
of Theorem 1, including even characteristic, without importing the
umbrella's C-group rank theorem.

For the remaining discussion \(q\) is odd and the rank is four. Every
adjacent product has order at least three. An order-one product would
repeat a marker. If an adjacent product has order two, divide the
marking at that edge. The nontrivial subgroups generated on its two
sides centralize one another and generate \(G\). Each is therefore
normal. Simplicity makes either one equal to \(G\), whereupon the other
lies in the trivial centre, a contradiction. This argument does not
assume that the two sides have trivial intersection in advance.

\subsection{4. Proper ends and the finite
interface}\label{proper-ends-and-the-finite-interface}

Let \(H=\langle r_0,r_1,r_2\rangle\) be proper in \(G=L(q)\) and put
\(D=\langle r_0,r_1\rangle\). The first adjacent order exceeds two, so
\(D\) is nonabelian. We apply the complete subgroup-type input of {[}E,
§2{]}.

A cyclic or elementary abelian group cannot contain \(D\). In \(A_4\)
all involutions lie in a four-group, so they cannot generate \(A_4\). In
an odd-characteristic affine group an involution has multiplier \(-1\)
and form \(x\mapsto-x+v\). Two such transformations commute precisely
when their translation parameters agree. The distinct commuting
endpoints \(r_0,r_2\) therefore exclude the affine alternative,
including its unipotent dihedral subgroups.

Suppose \(H\) is dihedral. Since \(r_0,r_2\) are distinct commuting
involutions, \(H\) has even rotation order and has a central involution
\(z\). Indeed, an odd dihedral group has only reflections, and two
reflections commute only when they are equal. The involution \(z\)
centralizes \(D\). So does \(r_3\). The common-centralizer bound forces
\(r_3=z\in H\), contradicting \(G=\langle H,r_3\rangle\) and properness.
This replaces both geometric and algebraic proofs of {[}LT, Lemma
3.9{]}; no assertion that every involution centralizer is maximal is
needed.

If \(H\) is a subfield group, the retained projective form of Dickson's
theorem conjugates it into \(P(k)\) for a proper subfield \(k\subsetneq
\mathbb F_q\). Theorem 5 of {[}E{]} places \(r_3\) in the same \(P(k)\),
since it centralizes \(D\). Hence \(G\subseteq P(k)\). Write \(s=|k|\).
Since \(q\ge s^2\),

\[ |G|=q(q^2-1)/2>s(s^2-1)=|P(k)|, \tag{2} \]

which is impossible. This handles both \(L(k)\) and \(P(k)\) subgroups,
including the small subfields; it uses their natural embeddings, not
only their abstract isomorphism types. We conclude that \(H\) is \(S_4\)
or \(A_5\). Reversing the marking proves the same assertion for \(G_0\).
This proves {[}LT, Theorem 3.10{]}.

The source's odd-end-label lemma also follows without IP. If
\(|r_0r_1|\) were even, its half-power would be the central involution
of \(D\). The common-centralizer bound would identify it with \(r_3\),
again putting \(r_3\) in the proper group \(H\). Thus the first label is
odd whenever \(H\) is proper. This recovers {[}LT, Lemma 3.5{]} with its
actual hypothesis. It does not assert that a marking with both ends
equal to \(G\) has odd end labels.

\textbf{Finite interface.} Every rank-three string marking of \(S_4\) or
\(A_5\) has IP. Up to automorphisms of the marked group its adjacent
orders and Petrie order are exactly

\[\begin{array}{c|ccc}
 S_4&(3,3,4)&(3,4,3)&(4,3,3)\\
 A_5&(3,5,5)&(5,3,5)&(5,5,3).
\end{array}\tag{3}\]

Here the third entry is \(|r_0r_1r_2|\). The new exhaustive program
\texttt{phase258\_local\_sggi.g} enumerates all ordered triples of
distinct involutions with commuting endpoints and tests generation. IP
is an output test, never an admission filter. There are 72 such triples
in \(S_4\) and 360 in \(A_5\); all pass the 64 subset-pair
intersections. Full automorphism orbits give (3). An explicit dihedral
string marking failing IP is a control against extending this assertion
to arbitrary groups. The exact presentations for (3) and the ten
compatible rank-four universals are already proved by {[}E, Lemmas
10--11{]} and its single bound GAP program. We reuse that calculation in
full.

Suppose both end groups are proper. The local census now licenses the
ten universal presentations, even though global IP is still unknown.
Every possible \(G\) is a quotient preserving the four generators. Four
rows are trivial. In the rows of orders 60 and 24 a stipulated end group
is already the whole universal group, and remains the whole image in
every quotient; this contradicts the required proper end. The order-96
row has no quotient \(L(q)\) in odd characteristic: \(60\) does not
divide 96, and for \(q\ge7\) the order is at least 168. The order-120
row is \(S_5\), whose only nontrivial proper quotient is \(C_2\);
neither this nor \(S_5\) is \(L(q)\) for odd \(q\ge5\).

The remaining rows are the simple groups \(L(11)\) and \(L(19)\). Their
only nontrivial quotients are themselves. The increasing odd-field order
formula forces \(q=11\) or \(19\), and the bound universal markings
already satisfy IP. Their reversal automorphisms give uniqueness under
duality as well as ordered-marking uniqueness. This is the complete
proof of {[}LT, Proposition 4.2{]}, with the missing sggi-to-C interface
made explicit. The standard local-map identifications give
hemi-icosahedral facets for the 11-cell and hemi-dodecahedral facets for
the 57-cell; their vertex figures are the dual local maps. The exact
face counts, self-duality and markings are retained from the same
complete table.

\subsection{5. The half-turn extension
theorem}\label{the-half-turn-extension-theorem}

Let \(q\) be odd and let \(N=(q-\alpha)/2\) for \(q\equiv\alpha\pmod4\).
Then \(N\) is even. A split or nonsplit cyclic torus of \(L(q)\) of
order \(N\) has a dihedral normalizer of order \(2N\). For clarity, the
split torus comes from diagonal projectivities of square eigenvalue
ratio and has order \((q-1)/2\). The nonsplit torus comes from
multiplication on \(\mathbb F_{q^2}\) modulo scalar multiplications,
followed by the determinant-square condition, and has order \((q+1)/2\).
Inversion of the two eigenlines supplies the other normalizer coset.
These standard torus facts do not assert maximality of the normalizer.

Suppose \((r_0,r_1,r_2)\) generates \(L(q)\), with distinct involutions,
commuting endpoints, and \(|r_0r_1|=k\) an odd divisor of \(N\).
Necessarily \(k\ge3\). The element \(r_0r_1\) is semisimple and has two
eigenlines over at most a quadratic field extension. The two possible
PSL torus orders \((q-1)/2\) and \((q+1)/2\) are coprime; since \(k>1\)
divides \(N\), its torus has order \(N\). Since
\(r_0(r_0r_1)r_0=(r_0r_1)^{-1}\), both reflections preserve this
unordered eigenline pair. Thus \(D=\langle r_0,r_1\rangle\) lies in the
associated dihedral torus normalizer. Its cyclic torus has a unique
involution \(z\), which centralizes the normalizer and hence \(D\).

The involution \(z\) is not in \(D\): that dihedral group has odd
rotation order and trivial centre, whereas \(z\) centralizes it. In
particular \(z\ne r_0,r_1\). It also differs from \(r_2\), since \(z\)
commutes with \(r_1\) whereas \(|r_1r_2|\ge3\). Therefore
\((r_0,r_1,r_2,z)\) is a string marking of rank four, and still
generates the original group. By the group--maniplex correspondence its
facets are the original map: there is one facet because the first three
markers already generate \(G\). This proves Theorem 2(1) and the
complete scope of {[}LT, Theorem 4.4{]}.

There is a useful local version. If a proper maximal subgroup \(H\le G\)
has a rank-three string marking, its first adjacent order is an odd
divisor of \(N\), and \(C_H(D)=1\), the same torus involution \(z\) lies
outside \(H\). Maximality gives \(\langle H,z\rangle=G\). The resulting
rank-four marking has facet group \(H\). The condition on the internal
centralizer is essential: membership in the ambient torus alone does not
prove that a new marker leaves the facet group.

\subsection{6. The explicit family with two full
ends}\label{the-explicit-family-with-two-full-ends}

For a prime \(p\equiv1\pmod{12}\), consider the projective classes of

\[R_0=\begin{pmatrix}0&1\\1&0\end{pmatrix},\quad
R_1=\begin{pmatrix}1&0\\1&-1\end{pmatrix},\quad
R_2=\begin{pmatrix}-1&0\\0&1\end{pmatrix},\quad
R_3=\begin{pmatrix}1&-2\\2&-1\end{pmatrix}.\tag{4}\]

The first three determinants are \(-1\) and the last is \(3\). Both are
squares modulo these primes: \(-1\) is a square since \(4\mid p-1\); if
\(\zeta\) has order three, then \((2\zeta+1)^2=-3\), and multiplication
by a square root of \(-1\) gives a square root of \(3\). Thus all four
classes lie in \(L(p)\). Their traces are zero and their determinants
nonzero, giving nonidentity involutions. Direct multiplication shows
that \(R_0,R_2\) commute projectively, that \((R_0R_1)^3\) is scalar and
\(R_0R_1\) is nonscalar, and that

\[[R_1R_2]=\left[\begin{pmatrix}1&0\\1&1\end{pmatrix}\right].\tag{5}\]

The latter element has exact order \(p\). Its powers give all lower
unipotent elementary matrices over \(\mathbb F_p\). Conjugation by
\(R_0\) gives all upper ones. These generate \(\operatorname{SL}_2(p)\):
left multiplication by lower and upper elementary matrices performs
elementary row elimination. Explicitly, writing \(u(t)\) and \(v(t)\)
for the upper and lower elementary matrices, put

\[h(t)=u(t)v(-t^{-1})u(t)=\begin{pmatrix}0&t\\-t^{-1}&0\end{pmatrix}.
\qquad h(t)h(-1)=\operatorname{diag}(t,t^{-1}).\]

The matrix \(h(1)\) interchanges the two coordinate lines, and the
displayed identity handles the last diagonal matrix after elimination.
Passing to projective classes proves
\(\langle r_0,r_1,r_2\rangle=L(p)\). This direct argument replaces a
subgroup-type exclusion for this step.

For the fourth matrix, set \(A=R_0R_1\). Then
\(2A-\operatorname{tr}(A)I=R_3\). Formula (1) shows that \(r_3\)
centralizes \(r_0,r_1\). The first product has odd order three, so
\(r_3\) lies outside \(\langle r_0,r_1\rangle\). It differs from \(r_2\)
because the latter has product order \(p\) with \(r_1\). Thus (4) is the
extension of §5. The opposite end contains \(r_1r_2\) of order
\(p\ge13\). By Theorem 1 it is \(S_4,A_5\) or \(L(p)\); the first two
groups have no element of that order. Both end groups are therefore
\(L(p)\), and the maniplex has one facet and one vertex.

These are the actual source generators, up to an explicit change of
basis. Choose \(a\in\mathbb F_p\) with \(a^2=-1\) and put
\(C=\operatorname{diag}(-a,1)\). If \(S_0,S_1,S_2\) denote the three
matrices printed in {[}LT, Lemma 4.5{]}, then the projective classes of
\(CS_iC^{-1}\) are exactly the first three classes in (4). Indeed the
three actual conjugated matrices are \(-aR_0,aR_1,aR_2\). This
identifies the construction at the level of its marked group, not just
by order or type. The fourth generator is the unique nonidentity common
centralizer element, so it agrees with the source's half-turn extension
as well.

The prime progressions are nonempty and infinite by Dirichlet's theorem.
Distinct primes give different group orders. Simplicity and §2 give
nonorientability. This proves Theorem 2(2) and {[}LT, Corollary 4.6{]}.

\subsection{7. Alternating and symmetric proper
ends}\label{alternating-and-symmetric-proper-ends}

Let \(p\equiv\pm1\pmod{20}\) be prime. Dickson's theorem gives a maximal
\(A_5\) subgroup \(H\) of \(L(p)\) {[}G, Theorem 2.2(6){]}. Give it the
local marking \((5,3,5)\) from (3), corresponding to the
hemi-dodecahedral map of type \(\{5,3\}\). Its first pair generates
\(D_{10}\). The internal centralizer \(C_{A_5}(D_{10})\) is trivial: the
centralizer of a 5-cycle in \(A_5\) is its cyclic group of order five,
and a reflection inverts that cycle, fixing only the identity in this
centralizer.

For these primes, the torus order \(N=(p-\alpha)/2\) is divisible by
ten, where \(p\equiv\alpha\pmod4\). The construction of §5 gives a new
involution outside \(H\), hence a rank-four marking generating \(L(p)\).
If \(p\ne19\), then \(p\) is neither 11 nor 19. The two-proper-end
classification in §4 forces the other end to be \(L(p)\). Thus this is
Class 2, with one vertex and \(|L(p)|/60\) hemi-dodecahedral facets. For
\(p=19\), the construction still gives a string marking, but the
two-proper-end theorem no longer forces the opposite end to be full. The
source excludes this value from its Class 2 assertion; no additional
classification at this value is needed here. Infinitely many admissible
primes remain, proving Theorem 2(3), {[}LT, Proposition 4.7{]} and
Corollary 4.8.

Now let \(p\equiv\pm1\pmod{24}\) be prime. Then \(p\ge23\), and
Dickson's theorem gives a maximal \(S_4\) subgroup \(H\) of \(L(p)\)
{[}G, Theorem 2.2(8){]}. Use its local marking \((3,4,3)\) in (3). The
first dihedral group is \(S_3\), of order six. Its centralizer in
\(S_4\) is trivial: this \(S_3\) fixes one point in the natural
degree-four action and permutes the other three transitively, and an
element centralizing that action must fix all four points. The relevant
torus order \(N=(p-\alpha)/2\) is divisible by twelve. The construction
of §5 therefore adjoins a centralizing involution outside \(H\) and
generates \(L(p)\) by maximality. Again the other end cannot be proper,
since \(p\notin\{11,19\}\). This proves Theorem 3, including the
complete all-prime range and the facet count. Dirichlet gives infinitude
in both residue classes.

Consequently the correct classification of end-group sizes is:

{\def\LTcaptype{none} % do not increment counter
\begin{longtable}[]{@{}lll@{}}
\toprule\noalign{}
Case & End groups, up to reversal & Vertices and facets \\
\midrule\noalign{}
\endhead
\bottomrule\noalign{}
\endlastfoot
Both full & \(G,G\) & \(1,1\) \\
One proper & \(G,S_4\) & \(1,|G|/24\) \\
One proper & \(G,A_5\) & \(1,|G|/60\) \\
Both proper & the 11-cell or 57-cell & \(11,11\) or \(57,57\) \\
\end{longtable}
}

The middle two cases cannot be self-dual, since their vertex and facet
counts differ. The first case can be self-dual; an exact example
follows. We make no complete enumeration of its possible markings or
dualities.

\subsection{8. Literal source corrections and finite
certificates}\label{literal-source-corrections-and-finite-certificates}

The new \texttt{phase258\_local\_sggi.g} supplies the precise category
bridge used in §4. It also checks a counterexample to {[}LT, Lemma
3.3{]} as printed: in \(L(5)\cong A_5\), an involution centralizer has
order four and lies properly in a subgroup \(A_4\) of order twelve. It
is not maximal in \(A_5\). The lemma also needs an odd-characteristic
qualification for its displayed dihedral formula. Our proof uses the
correct common centralizer theorem and torus normalizer, not this
maximality assertion. The broad uniqueness clauses of source Lemma 3.4
are not retained inputs.

For the omitted symmetric-end case, the certificate gives four explicit
permutations in the natural degree-24 action of \(L(23)\). They are
involutions, pairwise distinct, satisfy all string commutations and
generate the group of order 6072. Their adjacent orders are \(3,4,12\);
the first end has order 24 and is \(S_4\), while the other end is the
full group. The maniplex therefore has one vertex and 253 facets. These
facts directly contradict the introduction's restricted alternatives,
without relying on the all-prime proof or any conjectural
classification.

For the self-duality sentence on source page 9, use the following
marking of \(L(13)\) in degree 14:

\[\begin{aligned}
r_0&=(3\ 9)(4\ 10)(5\ 11)(6\ 12)(7\ 13)(8\ 14),\\
r_1&=(1\ 4)(2\ 14)(3\ 9)(5\ 13)(6\ 11)(7\ 12),\\
r_2&=(1\ 2)(3\ 7)(4\ 6)(8\ 14)(9\ 13)(10\ 12),\\
r_3&=(1\ 2)(4\ 14)(5\ 13)(6\ 12)(7\ 11)(8\ 10).
\end{aligned}\tag{6}\]

Both end groups have order 1092, and its type is \(\{3,7,3\}\). The
permutation

\[d=(1\ 3)(2\ 9)(4\ 7)(5\ 8)(6\ 12)(10\ 11)(13\ 14)\tag{7}\]

conjugates \(r_i\) to \(r_{3-i}\) for all four indices. Thus it
normalizes the generated group and induces the required reversal
automorphism. The resulting Class 1 maniplex is self-dual. This is an
explicit counterexample to the source's assertion about Class 1; the
elementary face-count exclusion applies to Class 2 instead. A
deliberately altered conjugator fails the four conjugacy identities in
the checker.

The source's paragraph following Proposition 4.7 first correctly calls
the local map of type \(\{5,3\}\) hemi-dodecahedral, then calls its
facets hemi-icosahedral. Our statement preserves the actual local
marking and uses the consistent hemi-dodecahedral name. None of these
corrections is presented as an answer to an open conjecture, and no
claim about unexamined editions is made.

\subsection{9. Remaining source scope, formal coverage and
cost}\label{remaining-source-scope-formal-coverage-and-cost}

The unnumbered Suzuki observation in {[}LT, §5{]} uses a separate
retained group fact: in the natural action of a simple Suzuki group,
every involution fixes exactly one point, and commuting involutions have
the same fixed point. In a rank-four string marking, \(r_0,r_2,r_3\)
would fix the same point, and then \(r_1\), which commutes with \(r_3\),
would fix it too. The generated group would fix a point, contrary to the
natural transitive action. The same argument applies to any higher rank,
since every later generator commutes with \(r_0\). This proves the
stated rank bound from that retained Suzuki input. The general abstract
centralizer version is already in the proof library; no new Suzuki
classification is claimed here. The nonsimple small group
\(\operatorname{Sz}(2)\) cannot be generated by involutions at all,
since its quotient \(C_4\) cannot be so generated.

The PGL observations and final four-family conjecture of source §5 are
kept separate from this PSL principal contract. In particular the
literal claim there about every proper nonabelian subgroup of \(P(q)\)
omits its affine, dihedral and subfield subgroups. It cannot serve as a
background classification. The companion work on that question retains
the actual maximal-subgroup theorem instead.

The source rank restriction, Theorem 3.10, Proposition 4.2, Theorem 4.4,
Lemma 4.5, Proposition 4.7 and Corollaries 4.6/4.8 are recovered in
full. The local maps and 11-/57-cell classification remain attributed to
their original authors and the reviewed common proof. The source
matrices are matched by an explicit conjugacy, the two infinite families
are proved for their full stated parameter ranges, and all necessary
small fields are handled. The omitted symmetric-end family and the
literal counterexamples are separate corrections. No complete census of
all rank-four PSL markings is claimed or needed by the source contract.

The new Lean file proves three general statements: commutation extends
from generators to the group; the first adjacent pair of a centre-free
string marking does not commute; and the common-centralizer uniqueness
condition forces rank at most four. These are the actual sggi bridge,
not a theorem assuming IP. The characteristic-two bridge and the
normal-form matrix calculations already have their separate, scoped
certificates. The projective conjugacy/descent, Dickson classification,
torus existence, finite universal orders, Dirichlet theorem and all
positive construction claims remain written or computational proof
obligations. They are not asserted to be covered by those three Lean
declarations.

The full {[}E{]} umbrella and finite table are required mathematical
inputs. The application adds the category bridge, complete
source-specific constructions and explicit corrections. Its local and
shared proof costs must be measured under the same convention as the
source; importing a previously proved theorem does not erase its
acquisition cost. This packet reports no compression ratio or
chronological sawtooth, and awards no new principal credit before
independent review.

\subsection{References}\label{references}

{[}LT{]} D. Leemans and M. Toledo, \emph{Maniplexes with automorphism
group \(\operatorname{PSL}(2,q)\)}, author manuscript dated 15 May 2023,
11 pages. Bound primary: \texttt{sources/psl2\_maniplex2023.pdf}.

{[}E{]} \emph{Centralizers, facets and finite amalgams on the projective
line}, frozen \texttt{phase238\_topic\_E\_umbrella.md}, with the
complete finite proof program \texttt{phase230\_finite\_check.g}. The
subgroup inputs and older classifications retain the attributions given
there.

{[}G{]} M. Giudici, \emph{Maximal subgroups of almost simple groups with
socle \(\operatorname{PSL}(2,q)\)}, arXiv:math/0703685v1 (23 March
2007), Theorem 2.2, physical and printed page 4. Bound primary:
\texttt{sources/phase257\_Giudici.pdf}. The maximal-subgroup
classification is retained, not recovered as an additional source here.

\end{document}
